\section{State formation and backend transfer boundaries}
\label{sec:transfer}
The B570 state experiments address a different question from the CUDA local repair: can attention-local control suffice when upstream projection and retained state depend on how prefill is partitioned? They hold continuation at a single query while intervening on prefill calls and the state consumed by later decoding.

On two exposed affected tasks, an external 128-token partition reproduces the native frontend's internally partitioned execution, including full initial FP32 logits and source-defined sequence KV, while the unpartitioned call differs at later output positions. Two longer controls have different numerical states but the same endpoint. A frozen partition sweep also finds a 129-token prefix counterexample, but the observations do not establish a universal 128-token threshold. Different state does not always change a greedy decision; matching endpoints do not prove that internal states are interchangeable.

Component-cross experiments reinforce this distinction. A continuation suffix can depend on retained KV under a qualified receiver contract, but an initial-logit swap can be a negative control when both sources select the same first token. Causal claims are limited to the tested donor/recipient states and exposed inputs.

Backend support is another independent boundary. The earlier frozen Vulkan operator rejected the lifted query axis because its support predicate required equal batch extents. The new scoped repair admits a narrow domain, but the first GQA adaptation changed AV evaluation and failed exact reference. Preserving the original AV reduction path restored byte equality on the registered synthetic panel and enabled a real q=2 request bridge. This sequence illustrates why support admission, operator correspondence and request qualification must be reported separately. The repair does not imply general support for arbitrary GQA strides or model features.
